\documentclass[11pt]{article}
\usepackage[margin=1in]{geometry}
\usepackage{amsmath,amssymb}
\usepackage{enumitem}

\newcommand{\Om}{\Omega_m}
\newcommand{\Ob}{\Omega_b}
\newcommand{\As}{A_s}
\newcommand{\ns}{n_s}
\newcommand{\Dhub}{D_H}
\newcommand{\xipm}{\xi_\pm}

\title{An Analytic Scaling for Cosmic-Shear\\
Data-Vector Emulation}
\author{V.\ Miranda}
\date{\today}

\begin{document}
\maketitle

\section{Idea}

A neural emulator learns the map from cosmological parameters
$\boldsymbol{\theta}_{\rm c}$ to the cosmic-shear data vector
$\mathbf d(\boldsymbol{\theta}_{\rm c})$. The dynamic range of
$\mathbf d$ across the prior is large, which makes the map hard
to fit. We remove most of that range \emph{analytically}: divide
out a fast approximate model and emulate the (flatter) residual.

Concretely, pick a reference (``mid'') cosmology and preprocess
each data vector element-wise by
\begin{equation}
  \tilde{\mathbf d}(\boldsymbol{\theta}_{\rm c})
  = \mathbf d(\boldsymbol{\theta}_{\rm c}) \odot
    \mathbf R(\boldsymbol{\theta}_{\rm c}),
  \qquad
  R \;=\; \frac{\xi^{\rm an}(\boldsymbol{\theta}_{\rm mid})}
               {\xi^{\rm an}(\boldsymbol{\theta}_{\rm c})},
  \label{eq:idea}
\end{equation}
where $\xi^{\rm an}$ is the analytic model below. The network
emulates $\tilde{\mathbf d}$ (which clusters near the mid
cosmology); the physical vector is recovered as
$\mathbf d=\tilde{\mathbf d}\oslash\mathbf R$ before the
$\chi^2$, so the data covariance and the reported metric are
unchanged. Because $R$ is a \emph{ratio}, any cosmology-common
factor (e.g.\ the bulk of the nonlinear boost) cancels: the
model only has to capture the cosmology \emph{dependence}.

\section{Approximations}

\begin{enumerate}[nosep]
  \item Limber projection.
  \item Each tomographic source bin is a single source plane at
        the peak of its $n(z)$ (a $\delta$-function $n(z)$).
  \item Pure dark energy expansion, $H(z)=H_0$, giving
        closed-form distances.
  \item Linear matter power spectrum (no halofit).
  \item Eisenstein--Hu zero-baryon transfer function.
  \item Hankel transform to real space approximated by a
        $\delta$-function at $\ell\theta=1$; the Limber integral
        evaluated at the lensing-kernel peak $u_\star$.
\end{enumerate}

None of these need to be accurate: $\xi^{\rm an}$ is only a
reference whose cosmology-dependence we divide out.

\section{Projection}

For a tomographic source pair $(i,j)$ the Limber shear spectrum is
\begin{equation}
  C_\ell^{ij}
  = \int_0^{\chi_{\min}}\! \frac{q_i(\chi)\,q_j(\chi)}{\chi^2}\,
    P\!\left(k=\tfrac{\ell+1/2}{\chi},\,z(\chi)\right)\,d\chi,
  \qquad \chi_{\min}=\min(\chi_i,\chi_j),
  \label{eq:limber}
\end{equation}
with the single-plane lensing efficiency
\begin{equation}
  q_i(\chi)=\frac{3}{2}\,\Om\left(\frac{H_0}{c}\right)^2
            \frac{\chi}{a(\chi)}\,\frac{\chi_i-\chi}{\chi_i},
            \qquad \chi<\chi_i .
  \label{eq:kernel}
\end{equation}
Under $H(z)=H_0$ the distance is closed form,
\begin{equation}
  \chi(z)=\frac{c}{H_0}\,z \equiv \Dhub\,z,
  \qquad a(z)=\frac{1}{1+z},
  \qquad \chi_i=\Dhub\,z_{s,i}.
  \label{eq:dist}
\end{equation}

\section{Linear power spectrum}

The matter spectrum follows from the primordial curvature spectrum
and the linear transfer. Write the dimensionless primordial spectrum
as
\begin{equation}
  \Delta_\mathcal{R}^2(k)=\As\,(k/k_p)^{\ns-1},
\end{equation}
with $k_p$ the pivot. In the matter era the sub-horizon density
contrast is, through the Poisson equation, proportional to the
primordial mode $\mathcal{R}(k)$,
\begin{equation}
  \delta_m(k,a)=\frac{2}{5}\,\frac{c^2 k^2}{\Om H_0^2}\,
  T(k)\,D(a)\,\mathcal{R}(k),
\end{equation}
where $T(k)$ is the Eisenstein--Hu zero-baryon transfer function
(Appendix~\ref{app:eh}) and $D(a)$ the growth factor. Squaring and
converting to the dimensionful spectrum, $P=(2\pi^2/k^3)\Delta^2$,
\begin{equation}
  P_{\rm lin}(k,a)=\frac{2\pi^2}{k^3}
  \left[\frac{2}{5}\,\frac{c^2 k^2}{\Om H_0^2}\right]^{\!2}
  T^2(k)\,D^2(a)\,\Delta_\mathcal{R}^2(k).
\end{equation}
With growth frozen ($H=H_0$ is de~Sitter, so $D=\text{const}$,
absorbed into the normalization) and at $z=0$, collecting the powers
of $k$ ($k^{-3}\cdot k^4\cdot k^{\ns-1}=k^{\ns}$) gives
\begin{equation}
  P_{\rm lin}(k)=\frac{8\pi^2}{25}\,\As\,
  \left(\frac{c^2}{\Om H_0^2}\right)^{\!2}
  k_p^{1-\ns}\,k^{\ns}\,T^2(k).
  \label{eq:plin}
\end{equation}

\section{Amplitude cancellation}

The lensing kernel and the power spectrum carry opposite powers of
the amplitude parameters. From Eqs.~\eqref{eq:kernel}
and~\eqref{eq:plin},
\begin{equation}
  q_i q_j \propto \Om^2\left(\frac{H_0}{c}\right)^{\!4},
  \qquad
  P_{\rm lin}\propto \As\,c^4\,\Om^{-2} H_0^{-4}\,k^{\ns}T^2(k),
\end{equation}
so in the Limber integrand the amplitude factors cancel term by
term,
\begin{equation}
  q_i q_j\,P_{\rm lin}\propto
  \Om^2\left(\frac{H_0}{c}\right)^{\!4}\!\cdot
  \As\,c^4\,\Om^{-2} H_0^{-4}\,k^{\ns}T^2(k)
  = \As\,k^{\ns}T^2(k).
\end{equation}
Hence $C_\ell$ is proportional to $\As$ alone: $\Om$ and $H_0$
cancel out of the \emph{amplitude},
\begin{equation}
  \boxed{\,C_\ell = \As\;G(\ell;\Om,H_0,\ns,z_s)\,}.
  \label{eq:Ascancel}
\end{equation}
In the $(\As,\Om,h)$ parametrization, $\As$ therefore carries the
entire lensing amplitude; $\Om$ and $H_0$ enter only through the
\emph{shape} $G$ (the transfer function and the geometry). This is
why dividing by $\As^{\rm mid}/\As$ is exact and removes the single
largest direction of variance.

\section{Real space and the effective wavenumber}

The correlation functions are Hankel transforms of $C_\ell$,
\begin{equation}
  \xipm(\theta)=\frac{1}{2\pi}\int_0^\infty \ell\,C_\ell\,
  J_{0,4}(\ell\theta)\,d\ell .
  \label{eq:hankel}
\end{equation}
Approximate the oscillating Bessel kernel by a $\delta$-function at
its argument, $J_{0,4}(\ell\theta)\approx a\,\delta(\ell\theta-1)$
for a constant $a$. Since
$\delta(\ell\theta-1)=\delta(\ell-1/\theta)/\theta$, the integral
collapses onto a single multipole $\ell_\star=1/\theta$,
\begin{equation}
  \xipm(\theta)\approx \frac{a}{2\pi}\int_0^\infty \ell\,C_\ell\,
  \frac{\delta(\ell-1/\theta)}{\theta}\,d\ell
  = \frac{a}{2\pi}\,\ell_\star^2\,C_{\ell_\star},
  \qquad \ell_\star=\frac{1}{\theta}.
  \label{eq:hankel2}
\end{equation}
The Limber integral~\eqref{eq:limber} is collapsed the same way, at
the lensing-kernel peak $\chi_\star=u_\star\chi_s$, fixing one
wavenumber
\begin{equation}
  k_\star=\frac{\ell_\star}{\chi_\star}
        =\frac{1}{\theta\,u_\star\chi_s}
        =\frac{H_0}{c\,u_\star\,\theta\,z_s},
\end{equation}
where the last step used $\chi_s=\Dhub z_s=(c/H_0)z_s$
from Eq.~\eqref{eq:dist}. The transfer function only ever sees the
dimensionless combination $q\equiv k_\star\Theta^2/(\Om h^2)$; with
$H_0=100\,h~{\rm km\,s^{-1}Mpc^{-1}}$ the factor
$H_0/(\Om h^2)=100/(\Om h)=100/\Gamma$, so
\begin{equation}
  \boxed{\;
  q \equiv \frac{k_\star\,\Theta^2}{\Om h^2}
    = \frac{K}{\theta\,z_s\,\Gamma},
  \qquad
  K=\frac{100\,\Theta^2}{c\,u_\star},
  \qquad
  \Gamma\equiv\Om h,
  \qquad
  \Theta\equiv\frac{T_{\rm CMB}}{2.7\,{\rm K}},
  \;}
  \label{eq:q}
\end{equation}
with $\theta$ in radians. The shape parameter $\Gamma=\Om h$ and
the source redshift enter \emph{only} through the single
combination $q\propto 1/(\theta\,\Gamma\,z_s)$: a rescaling of the
angle is equivalent to a change in $\Gamma z_s$.

\section{The scaling}

Put the pieces together. With the amplitude stripped
(Eq.~\ref{eq:Ascancel}), the Limber spectrum~\eqref{eq:limber} is
\begin{equation}
  C_\ell = \As\,c_0\!\int_0^{\chi_s}\!\! \frac{d\chi}{\chi^2}\,
  \left[\frac{\chi}{a}\,\frac{\chi_s-\chi}{\chi_s}\right]^2
  k^{\ns}T^2(k), \qquad k=\ell/\chi,
\end{equation}
with $c_0$ a constant. Substitute $u=\chi/\chi_s$, so $\chi=u\chi_s$,
$z=u z_s$, $a=1/(1+u z_s)$, and $d\chi=\Dhub z_s\,du$. The geometric
bracket simplifies,
\begin{equation}
  \frac{1}{\chi^2}
  \left[\frac{\chi}{a}\,\frac{\chi_s-\chi}{\chi_s}\right]^2
  =\frac{\big[u\chi_s(1+u z_s)(1-u)\big]^2}{(u\chi_s)^2}
  =(1+u z_s)^2(1-u)^2,
\end{equation}
so $C_\ell=\As\,c_0\,\Dhub z_s\!\int_0^1\!du\,(1+u z_s)^2(1-u)^2\,
k^{\ns}T^2(k)$. Collapse this integral at the kernel peak $u_\star$
and use $\xipm\propto\ell_\star^2 C_{\ell_\star}$ with
$\ell_\star=1/\theta$ (Eq.~\ref{eq:hankel2}):
\begin{equation}
  \xipm(\theta)\propto \As\,
  \Dhub z_s\,(1+u_\star z_s)^2(1-u_\star)^2\,
  \ell_\star^2\,k_\star^{\ns}\,T^2(k_\star).
  \label{eq:xicollapse}
\end{equation}
At fixed $\theta$ the factor $\ell_\star^2=\theta^{-2}$ and the
$u_\star$ constants are common to every cosmology and cancel in the
ratio $R$ below, so we keep only the cosmology-dependent factors.
Splitting $k_\star^{\ns}=q^{\ns}k_\Gamma^{\ns}$ with
$k_\Gamma=\Om h^2/\Theta^2$ (so $k_\star=q\,k_\Gamma$) and using
$\Dhub=c/H_0\propto1/h$,
\begin{equation}
  \xipm(\theta)=\As\;\mathcal{N}(\Om,h,z_s,\ns)\;
                q^{\ns}\,T^2(q),
  \qquad
  \mathcal{N}\propto
  \frac{z_s\,(1+u_\star z_s)^2\,(\Om h^2)^{\ns}}{h},
  \label{eq:xifull}
\end{equation}
where $\mathcal{N}$ is a slowly varying geometric amplitude and
$q^{\ns}T^2(q)$ is the dimensionless shape.

\section{The preprocessing ratio}

In the ratio, the survey-fixed parts of $\mathcal{N}$ ($z_s$ and the
geometric factor $(1+u_\star z_s)^2$) are identical for the reference
and the target and cancel exactly. What survives is a per-cosmology
\emph{scalar} amplitude $(\Om h^2)^{\ns}/h$: the $\Om$--$H_0$
amplitude that cancelled out of the leading $C_\ell$
(Eq.~\ref{eq:Ascancel}) re-enters here through the tilt
$k_\Gamma^{\ns}=(\Om h^2)^{\ns}$ and the $1/h$ distance measure. It
is a \emph{second} amplitude direction, independent of $A_s$, and is
kept. Equation~\eqref{eq:idea} becomes the per-element factor
\begin{equation}
  \boxed{\;
  R(\theta)=R_{\rm amp}\;
  \frac{q_{\rm mid}^{\,\ns^{\rm mid}}\,T^2(q_{\rm mid})}
       {q^{\,\ns}\,T^2(q)},
  \qquad
  q=\frac{K}{\theta\,z_{\rm eff}\,\Gamma},
  \;}
  \label{eq:R}
\end{equation}
with the $\theta$-independent amplitude factor
\begin{equation}
  R_{\rm amp}=\frac{\As^{\rm mid}}{\As}\;
  \frac{(\Om^{\rm mid}h_{\rm mid}^2)^{\ns^{\rm mid}}/h_{\rm mid}}
       {(\Om h^2)^{\ns}/h},
  \label{eq:Ramp}
\end{equation}
and $q_{\rm mid}=K/(\theta\,z_{\rm eff}\,\Gamma_{\rm mid})$. Each
cosmology uses its own $(\As,\ns,\Gamma,\Om,h)$ and the reference its
own; $R=1$ at the reference. The $\theta$-dependence enters only
through the non-power-law transfer ratio $T^2(q_{\rm mid})/T^2(q)$
(the genuine shape correction), since
$q_{\rm mid}/q=\Gamma/\Gamma_{\rm mid}$ is itself
$\theta$-independent.

\paragraph{Tomographic cross-pairs.}
The cross spectrum~\eqref{eq:limber} is cut off at the nearer
source, so its effective lensing scale is that of the lower-redshift
bin. The single-scale collapse therefore uses
\begin{equation}
  z_{\rm eff}=\min(z_{s,i},\,z_{s,j}),
  \label{eq:zeff}
\end{equation}
not the geometric mean; for auto-pairs ($i=j$) they coincide. Since
$z_{\rm eff}$ cancels at leading order in $R$, this choice only
shifts the (mild) transfer correction, mattering at the few-percent
level for asymmetric pairs.

\section{Empirical validation}

On a held-out validation set, measuring per element the spread ratio
across cosmologies $\sigma[\tilde d_e]/\sigma[d_e]$ (its square is
the fraction of variance \emph{remaining}) and the fraction of
elements whose spread the rescaling reduces:

\begin{center}
\begin{tabular}{lccc}
\hline
 & median ratio & mean ratio & frac.\ improved \\
\hline
$A_s$ only        & 0.79 & 0.76 & 1.00 \\
$+$ shape         & 0.52 & 0.60 & 0.89 \\
$+\,\mathcal{N}$  & 0.46 & 0.50 & 0.98 \\
\hline
\end{tabular}
\end{center}

\noindent The exact $A_s$ rescaling alone removes ${\sim}37\%$ of the
variance; the shape term roughly doubles that (${\sim}74\%$); the
$\mathcal{N}$ amplitude takes it to ${\sim}79\%$. Crucially
$\mathcal{N}$ also raises the improved fraction from $0.89$ to
$0.98$. The shape term alone leaves the $(\Om h^2)^{\ns}/h$ amplitude
in the residual, which dominates at large $\theta$, where $T\to1$
and the data variation is nearly pure amplitude---exactly the
elements the shape term made worse. Restoring $\mathcal{N}$ removes
that offset, and it is most reliable there (large scales, no
transfer curvature). The reduction is largest at small $\theta$ (the
well-constrained, small-error-bar scales) and falls to ${\sim}1$
only at the largest $\theta$, which carry the biggest covariance and
are down-weighted in the $\chi^2$.

\section{Validity and limitations}

\begin{itemize}[nosep]
  \item \textbf{Baryons.} The zero-baryon $T(k)$ has \emph{no}
        $\Ob$ dependence, so $R$ cannot remove the
        $\Ob h^2$ (physical baryon density) direction. That is
        precisely the direction that drives the hard,
        high-$\Ob h^2$ cosmologies, so this scaling flattens the
        broadband bulk but leaves the baryon tail to the network.
        Capturing it would require the wiggle (with-baryon)
        Eisenstein--Hu transfer.
  \item \textbf{Nonlinearity.} The model is linear, but $R$ is a
        ratio: the cosmology-common nonlinear boost largely
        cancels, so the rescaling remains accurate at small scales
        (where it empirically performs best) as long as that boost
        varies weakly across the prior.
  \item \textbf{Large scales.} At large $\theta$ the data variation
        is nearly pure amplitude ($T\to1$); the $\mathcal{N}$ scalar
        removes it, so these elements flatten once it is included
        (the improved fraction rises to $0.98$). The small residual
        sits in the highest-covariance bins and is down-weighted in
        the $\chi^2$.
\end{itemize}

\appendix
\section{Eisenstein--Hu zero-baryon transfer function}
\label{app:eh}

With $\Theta=T_{\rm CMB}/2.7\,{\rm K}$ and $k$ in
${\rm Mpc}^{-1}$,
\begin{align}
  q   &= \frac{k\,\Theta^2}{\Om h^2}, \\
  L   &= \ln\!\left(2e + 1.8\,q\right), \\
  C   &= 14.2 + \frac{731}{1+62.5\,q}, \\
  T(k)&= \frac{L}{L + C\,q^2}.
\end{align}

\end{document}
